\documentclass[a4paper]{article}

\usepackage[brazilian]{babel}
\usepackage[utf8]{inputenc}
\usepackage[T1]{fontenc}
\usepackage{lmodern}
\usepackage{amssymb}
\usepackage{amsmath}
\usepackage{graphicx}
\usepackage{subfig}
\usepackage{float}
\usepackage{framed}
\usepackage{xcolor}
\usepackage{colortbl}
\usepackage{subcaption}
\usepackage{listings}
\usepackage[most]{tcolorbox}

\newcommand{\todo}[1] {{\color{red} TODO:#1}}
\newcommand{\answer}[1] {{\color{blue} \textbf{Resposta}: #1}}

\usepackage{tikz} % Para o fillblank
% Para ter o tab e o fillblank
\newcommand*\circled[1]{\tikz[baseline=(char.base)]{
\node[shape=circle,draw,inner sep=0.5pt] (char) {#1};}}
\newcommand{\circleup}[1]{\textsuperscript{\circled{#1}}}
\newcommand{\fillblank}[2]{{\rule{#1cm}{.2mm}}\circleup{#2}}
\usetikzlibrary{arrows}

\tikzset{
treenode/.style = {align=center, inner sep=0pt, text centered,
font=\sffamily},
arn_n/.style = {treenode, circle, black, draw=black, text width=2.5em},
arn_r/.style = {treenode, circle, black, draw=black, 
text width=2.5em},  
arn_x/.style = {treenode, rectangle, draw=black,
minimum width=1em, minimum height=0.5em}
}

\newif\ifgabarito% Texto para aparecer apenas no gabarito
\gabaritofalse   % Texto para aparecer no modo prova

%-------------------------------
%\usepackage{pcsamc}
%\usepackage{comment}
%\includecomment{resolucao}
%\excludecomment{gabarito}
%\gabaritofalse
%-------------------------------


%Para o gabarito comentar o bloco acima e descomentar esse
%-------------------------------
\usepackage{comment}
\includecomment{gabarito}
\includecomment{resolucao}
\usepackage{pcsamc-gabarito}
\gabaritotrue
%-------------------------------

\newcommand{\tab}{\ \ \ \ }

\usepackage{array} % Para tabular
\newcolumntype{L}[1]{>{\raggedright\let\newline\\\arraybackslash\hspace{0pt}}m{#1}}
\newcolumntype{C}[1]{>{\centering\let\newline\\\arraybackslash\hspace{0pt}}m{#1}}
\newcolumntype{R}[1]{>{\raggedleft\let\newline\\\arraybackslash\hspace{0pt}}m{#1}}


% Para evitar substituição de << por aspas francesas
\usepackage{microtype}
\DisableLigatures[<]{encoding = T1}

% Estilo PCSAMC com as personalizações
\usepackage{enumitem}
\usepackage{xhfill}
\usepackage{multicol}
\setlength{\columnsep}{1.5cm}
\setlength{\columnseprule}{1pt}

\begin{document}

%##################################################################%
\element{testes}{
{
\def\AMCbeginQuestion#1#2{\par\noindent{\bf (POSCOMP 2004 - Questão 39)}}
\begin{question}{} Em um sistema operacional, um processo pode, em um dado instante de tempo, estar em um de três estados: em execução, pronto ou bloqueado. Considere as afirmativas abaixo sobre as possíveis transições entre estes estados que um processo pode realizar.
\begin{enumerate}[label=\Roman*.]
\item Do estado em \emph{execução} para o estado \emph{bloqueado}
\item Do estado em \emph{execução} para o estado \emph{pronto}
\item Do estado \emph{pronto} para o estado em \emph{execução}
\item Do estado \emph{pronto} para o estado \emph{bloqueado}
\item Do estado \emph{bloqueado} para o estado em \emph{execução}
\item Do estado \emph{bloqueado} para o estado \emph{pronto}
\end{enumerate}
\begin{multicols}{2}\columnseprule=0.1pt
\begin{choices}
\correctchoice{Somente as afirmativas I, II, III e VI são verdadeiras.}
\wrongchoice{Somente as afirmativas I, II e III são verdadeiras.}
\wrongchoice{Somente as afirmativas I, III, IV e VI são verdadeiras.}
\wrongchoice{Somente as afirmativas I, III, IV e V são verdadeiras.}
\wrongchoice{Todas as afirmativas são verdadeiras.} 
\end{choices}
\end{multicols}
\end{question}
}

{
\def\AMCbeginQuestion#1#2{\par\noindent{\bf (POSCOMP 2014 - Questão 50)}}
\begin{question}{} Em relação ao gerenciamento de processos, atribua V (verdadeiro) ou F (falso) às afirmativas a seguir.
\begin{enumerate}[label=(\ \ )]
\item Na espera ocupada, o processo é transferido para estado de bloqueado até que sua fatia de tempo termine e então ele retorna para fila de prontos.

\item O bloco de controle de processos (BCP -- \textit{Process Control Block}) é utilizado para armazenar informações sobre processos, e essas informações são utilizadas na troca de contexto de processos.

\item \textit{Threads} apresentam menor custo de criação quando comparadas aos processos, pois compartilham alguns elementos do processo, como espaço de endereçamento.

\item Um processo pode estar nos seguintes estados: pronto, aguardando execução, em execução e bloqueado.

\item Um processo pode ser criado por uma chamada de sistema \texttt{fork()}, nesse caso, o processo gerado (conhecido como ``filho'') é uma cópia exata do processo original, com os mesmos valores de variáveis em memória, diferenciando-se apenas no identificador do processo.
Assinale a alternativa que contém, de cima para baixo, a sequência correta.
\end{enumerate}

\begin{multicols}{3}\columnseprule=0.1pt
\begin{choices}
\correctchoice{F, V, V, F, V.}
\wrongchoice{V, V, F, V, F.}
\wrongchoice{V, F, V, F, F.}
\wrongchoice{V, F, F, F, V.}
\wrongchoice{F, F, F, V, V.} 
\end{choices}
\end{multicols}
\end{question}
}
%
%
%
{
\def\AMCbeginQuestion#1#2{\par\noindent{\bf (POSCOMP 2019 - Questão 45)}}
\begin{question}{} Considere o programa abaixo escrito em linguagem C. No instante da execução da linha 5, ter-se-á uma hierarquia composta de quantos processos e threads, respectivamente?
\begin{lstlisting}[language=C]
main(){
  int i;

  for(i=0;i<3;i++)
    fork();

  while(1);
}
\end{lstlisting}
\begin{multicols}{5}\columnseprule=0.1pt
\begin{choices}
\correctchoice{8 e 8.}
\wrongchoice{1 e 0.}
\wrongchoice{3 e 0.}
\wrongchoice{4 e 1.}
\wrongchoice{7 e 7.} 
\end{choices}
\end{multicols}
\end{question}
}
%
%
%
{
\def\AMCbeginQuestion#1#2{\par\noindent{\bf (POSCOMP 2022 - Questão 46)}}
\begin{question}{} O programa (em linguagem C) abaixo executa em um sistema operacional da família UNIX. Considere que todas as rotinas invocadas no programa executam sem erro. Assinale a alternativa que indica o resultado impresso na tela pelo programa.
\begin{lstlisting}[language=C]
signed int i;
int main(void){
  if (fork() > 0)
	i++;
  else
	i++;
  i++;
  printf("%d ", i);
}
\end{lstlisting}

\begin{multicols}{2}
\begin{choices}
\correctchoice{2 2}
\wrongchoice{1 1}
\wrongchoice{3 3}
\wrongchoice{4 4}
\wrongchoice{Indeterminado Indeterminado} 
\end{choices}
\end{multicols}
\end{question}
}
%
%
%
{
\def\AMCbeginQuestion#1#2{\par\noindent{\bf (POSCOMP 2025 - Questão 42)}}
\begin{question}{} Tipicamente, a execução de um processo em um sistema operacional de tempo compartilhado caracteriza-se pela alternância entre ciclos de execução da CPU e ciclos de espera de E/S, até o seu término. Os tempos médios dos ciclos de execução da CPU e de espera de E/S denominaremos, respectivamente, exec\_t e es\_t. Considere um conjunto de processos que seguem esse padrão de alternância entre CPU e E/S, os quais executam sob um sistema operacional com escalonamento do tipo round-robin (RR) preemptivo. Assuma que o tempo médio para realizar uma troca de contexto seja dado por switch\_t. Assinale a alternativa que melhor minimiza o tempo médio de turnaround dos processos.
%turnaround = tempo de término − tempo de chegada
	

\begin{multicols}{2}
\begin{choices}
\correctchoice{(exec\_t > switch\_t) e (exec\_t < quantum).}
\wrongchoice{(exec\_t > switch\_t) e (exec\_t > quantum).}
\wrongchoice{(exec\_t < switch\_t) e (exec\_t < quantum).}
\wrongchoice{(exec\_t < switch\_t) e (exec\_t > quantum).}
\wrongchoice{(exec\_t + switch\_t) > quantum.} 
\end{choices}
\end{multicols}
\end{question}
}
%
%
%
{
\def\AMCbeginQuestion#1#2{\par\noindent{\bf (PCS3746 - 2025)}}
\begin{question}{} Considere o código em Assembly abaixo que utiliza \texttt{fork} para criar dois processo. Assuma que o código compila e executa sem gerar erros ou exceções em tempo de execução.

\begin{multicols}{2}\columnseprule=0.1pt
\begin{lstlisting}[
language={[x86masm]Assembler},
basicstyle=\small,
keywordstyle=\color{blue},
commentstyle=\color{gray},
stringstyle=\color{red}
]

section .data
char_A db 'A', 0
char_B db 'B', 0

section .text
global _start

_start:
;Intrucoes para realizar o fork()
mov eax, 2
int 0x80	

cmp eax, 0
jnz B ;Se EAX for != 0, salta para B:
jmp finalizar
;
;
;
;
;
;
\end{lstlisting}
\begin{lstlisting}[
language={[x86masm]Assembler},
basicstyle=\small,
keywordstyle=\color{blue},
commentstyle=\color{gray},
stringstyle=\color{red}
]
A:
; Instrucoes para mostrar char_A
mov ebx, 1
mov ecx, char_A
mov edx, 3
mov eax, 4
int 0x80

jmp finalizar

B:
; Instrucoes para mostrar char_B
mov ebx, 1
mov ecx, char_B
mov edx, 3
mov eax, 4
int 0x80

finalizar:
mov eax, 1
xor ebx, ebx
int 0x80
\end{lstlisting}
\end{multicols}

Assinale todas as alternativas \textbf{verdadeiras}.	
\begin{choices}
\correctchoice{Os processo pai e filho executam o trecho de código do rótulo \texttt{finalizar}.}\scoring{b=1,e=0,v=0}
\correctchoice{O único processo que executará o trecho de código do rótulo \texttt{B} será o processo pai.}\scoring{b=1,e=0,v=0}%Verdadeiro. O filho salta direto para finalizar e o pai executa porque está logo abaixo

\wrongchoice{A execução desse código exibirá os caracteres AB ou BA dependendo do escalonamento feito pelo sistema operacional.}%F. O rótulo A nunca executará
\wrongchoice{O processo filho executará o trecho de código do rótulo \texttt{A}.}%F O processo filho executará somente o trecho finalizar
\wrongchoice{O processo filho executará o trecho de código do rótulo \texttt{B}.}%F mesma questão que acima
\wrongchoice{Se o sistema operacional utilizar o algoritmo de escalonamento do tipo alternância circular (\emph{Round-Robin}), a execução do código exibirá os caracteres AB ou BA.}%F. Não importa o tipo de escalonamento, o fluxo do filho sempre irá direto para o rótulo de finalizar
\wrongchoice{Se o sistema operacional utilizar um algoritmo de escalonamento do tipo não-preemptivo, a execução do código exibirá os caracteres AB ou BA.}%F. Não importa o tipo de escalonamento, o fluxo do filho sempre irá direto para o rótulo de finalizr 
\end{choices}
\end{question}
}
%
%
%
{
\def\AMCbeginQuestion#1#2{\par\noindent{\bf (PCS3746 - 2025)} #2\hspace*{1em}}
\begin{question}{}
Assuma um cenário hipotético em que um sistema operacional \textbf{não garante} que os registradores de propósito geral são salvos e recuperados no início e término da rotina de tratamento de interrupções, respectivamente.
%garante que as flags de estado da CPU e todos os registadores \textbf{exceto} os registradores EAX, EBX, ECX e EDX são salvos e recuperados, respectivamente, no início e término da rotina de tratamento de interrupções.
Assinale a solução que deveria ser implementada para que o tratamento de interrupções funcione corretamente no sentido de garantir que a execução dos processos não seja afetada.

\begin{choices}
\correctchoice{Utilizar as instruções {\tt pusha} e {\tt popa} no início e fim (antes de  {\tt iret}), respectivamente, da(s) rotina(s) de tratamento de interrupções.}
%\wrongchoice{Utilizar somente interrupções de software.}%Interrupção é interrupção, de qualquer forma vai ser chamado alguma coisa
%\wrongchoice{Acessar o registrador de apoio e aumentar seu valor para aumentar o intervalo entre interrupções}%Falso. Isso irá retardar o problema
\wrongchoice{Aumentar o intervalo em que as interrupções ocorrem.}
\wrongchoice{Utilizar a instrução {\tt xchg} para garantir atomicidade entre as operações dentro da rotina de tratamento de interrupções.}
\wrongchoice{Utilizar as instruções {\tt pushf} e {\tt pop} no início e fim (antes de  {\tt iret}), respectivamente, da(s) rotina(s) de tratamento de interrupções.}
\wrongchoice{Utilizar as instruções  {\tt cli} e {\tt sti} para desabilitar e habilitar interrupções enquanto o tratamento de uma interrupção é executado.}
\wrongchoice{Mover o registrador de segmento de código (\texttt{CS}) para outra posição de memória.}
\end{choices}
\end{question}
}
%
%
%
{
\def\AMCbeginQuestion#1#2{\par\noindent{\bf (PCS3746 - 2025)} #2\hspace*{1em}}
\begin{question}{} No exemplo visto em aula de interrupções no DoS (utilizando o DosBox), o endereço de tratamento de interrupções do tipo \texttt{int 21}\textbf{h} estava localizado na posição de memória $0\times0084$, onde os dois primeiros bytes indicavam o endereço da rotina e os dois últimos bytes o segmento de código. 
%$0\times0084$  = $0\times21$ * $0\times4$
Suponha que um sistema operacional implementará uma interrupção denominada \texttt{int 4}\textbf{h}. Seguindo o mesmo raciocínio adotado para a \texttt{int 21}\textbf{h}, para que o sistema operacional saiba quais instruções executar ao tratar uma interrupção do tipo \texttt{int 4}\textbf{h}, em qual posição de memória ser armazenado o endereço da rotina que trata a interrupção \texttt{int 4h}? Observação importante: nas alternativas desta questão ignore aspectos relacionados ao segmento de código (\texttt{CS}) e o formato de armazenamento \emph{little-endian}.
%Em qual posição de memória deve ser gravado o endereço da rotina que será executada quando ocorrer a interrupção \texttt{int 4}?

\begin{multicols}{5}\columnseprule=0.1pt
\begin{choices}
\correctchoice{$0\times0010$}\scoring{b=1,e=0,v=0}
\wrongchoice{$0\times0016$}
\wrongchoice{$0\times0000$}
\wrongchoice{$0\times 0084$}
\wrongchoice{$0\times0085$}
\end{choices}
\end{multicols}

Considere a representação conceitual de endereço de memória $\times$ conteúdo (instrução) abaixo. Nesta representação, todos os endereços estão em \textbf{hexadecimal} (o prefixo $0\times$ foi omitido por simplicidade).
\begin{center}
{\ttfamily
\begin{tabular}{|c|c|c|c|c|c|c|c|c|c|c|l|c|c|}
\cline{1-2} \cline{4-5} \cline{7-8} \cline{10-11} \cline{13-14}
0106 & pusha &  & 010C & pusha &  & 0112 & pusha &  & 0118 & pusha &  & 011E & pusha \\ \cline{1-2} \cline{4-5} \cline{7-8} \cline{10-11} \cline{13-14} 
0107 & pushf &  & 010D & pushf &  & 0113 & pushf &  & 0119 & pushf &  & 011F & pushf \\ \cline{1-2} \cline{4-5} \cline{7-8} \cline{10-11} \cline{13-14} 
0108 & add esp, 8 &  & 010E & add esp, 4 &  & 0114 & sub esp, 8 &  & 011A & pop rax &  & 0120 & pop eax \\ \cline{1-2} \cline{4-5} \cline{7-8} \cline{10-11} \cline{13-14} 
0109 & add esp, 2 &  & 010F & nop &  & 0115 & sub esp, 2 &  & 011B & nop &  & 0121 & pop ebx \\ \cline{1-2} \cline{4-5} \cline{7-8} \cline{10-11} \cline{13-14} 
010A & popa &  & 0110 & popa &  & 0116 & popa &  & 011C & popa &  & 0122 & nop \\ \cline{1-2} \cline{4-5} \cline{7-8} \cline{10-11} \cline{13-14} 
010B & iret &  & 0111 & iret &  & 0117 & iret &  & 011D & iret &  & 0123 & iret \\ \cline{1-2} \cline{4-5} \cline{7-8} \cline{10-11} \cline{13-14} 
\end{tabular}%
}
\end{center}

A partir dessa memória e considerando uma \textbf{arquitetura de 32 bits}, qual o conteúdo deve ser colocado no endereço da \texttt{int 4}\textbf{h} (respondido na questão anterior) para que a rotina de tratamento desta interrupção tenha o seguinte comportamento: salvar registradores de propósito geral, salvar as flags de estado da CPU e recuperar o conteúdo dos registradores de propósito geral antes de retornar da interrupção.
%realize as seguintes operações, nesta ordem: (i) desabilitar interrupções, (ii) salvar registradores de propósito geral, (iii) empilhar o estado das flags de CPU, (iv) recuperar o conteúdo dos registradores de propósito geral, (v) habilitar interrupções e (vi) retornar da interrupção.
\begin{multicols}{5}
\begin{choices}
\correctchoice{$0\times010C$}
\wrongchoice{$0\times0106$}
\wrongchoice{$0\times0112$}
\wrongchoice{$0\times0118$}
\wrongchoice{$0\times011E$}
\end{choices}
\end{multicols}
\end{question}
}
%\input{Threads}
%\input{GerenciamentoMemoria}
%{

\def\AMCbeginQuestion#1#2{\par\noindent{\bf (POSCOMP 2024 – Questão 44)}}
\begin{question}{} Um computador precisa ter seu disco atual, formatado com uma única partição, substituído por outro disco. Antes de realizar a troca, foi realizada uma análise do tamanho dos
arquivos armazenados no disco atual, cujo resultado indicou que 98\% dos arquivos possuem tamanho
entre 1.000 bytes e 2.030 bytes (mediana de 1.515 bytes). Assinale a alternativa que lista os
parâmetros de formação do novo disco e que resulte no menor desperdício de espaço físico e menor
probabilidade de fragmentação do disco.
\begin{choices}
\correctchoice{Tamanho de bloco = 2.048 bytes e sistema de arquivos baseado em alocação encadeada.}
\wrongchoice{Tamanho de bloco = 1.024 bytes e sistema de arquivos baseado em alocação encadeada.}
\wrongchoice{Tamanho de bloco = 1.024 bytes e sistema de arquivos baseado em alocação indexada.}
\wrongchoice{Tamanho de bloco = 2.048 bytes e sistema de arquivos baseado em alocação contígua.} 
\wrongchoice{Tamanho de bloco = 1.024 bytes e sistema de arquivos baseado em alocação contígua.}
\end{choices}


\ifgabarito
{\color{red} Resposta ao recurso: MANTIDA alternativa 'D'. A questão apresenta o cenário de formatação de uma partição de disco, e solicita que “assinale a alternativa que lista os parâmetros de formação do novo disco e que resulte no menor desperdício de espaço físico e menor probabilidade de fragmentação do disco”. As opções disponíveis apresentam dois tamanhos de bloco (1024 bytes e 2048 bytes) e três métodos de alocação de espaço em disco (contígua, encadeada e indexada). O enunciado informa que 98\% dos arquivos a serem armazenados na nova partição possuem tamanho entre 1.000 bytes e 2030 bytes (mediana de 1.515 bytes). Importante ressaltar que nada é dito sobre a distribuição dos tamanhos dos arquivos. Em termos de método de alocação de espaço em disco, entre as três opções disponíveis, a literatura indica que o método que oferece a menor probabilidade de fragmentação de disco é a alocação encadeada, dada a sua flexibilidade em alocar blocos não contíguos e por não reservar espaço para tabela de índices. Desse modo, as alternativas disponíveis com esse método são as letras (A) e (D). A diferença entre essas duas alternativas está no tamanho de bloco. Em termos do tamanho de bloco, dada a ausência de informações sobre a distribuição dos tamanhos dos arquivos, nada pode ser inferido a esse respeito, além do seu valor mediano de 1.515 bytes que é informado. Neste sentido, para o tamanho de bloco de 1024 bytes (alternativa A), com base no valor mediano informado, um arquivo ocuparia dois blocos de 1024 bytes, sendo que a ocupação do segundo bloco seria parcial. Já para o tamanho de bloco de 2048 bytes (alternativa D), com base no valor mediano, um arquivo ocuparia um único bloco, também de forma parcial. O espaço não utilizado (desperdiçado) nas duas configurações de tamanho de bloco é o mesmo, 533 bytes. Portanto, em termos de desperdício, ambas as alternativas teriam o mesmo efeito, tendo como base o valor mediano informado. Contudo, em termos de fragmentação de disco, a opção (A), com dois blocos de 1024 bytes, tem uma maior probabilidade de fragmentação de disco, praticamente o dobro da opção (D), já que esses arquivos, ao serem removidos, deixariam dois blocos livres enquanto sua contraparte deixaria apenas um. Note que essa probabilidade de fragmentação de disco é acentuada pela alocação encadeada, haja vista que os dois blocos em questão não, necessariamente, estariam localizados de forma contígua. Em especial sobre o problema de fragmentação de disco, vale destacar que o enunciado não busca pela solução ideal para o problema, mas sim aquela que resulte na “menor probabilidade de fragmentação de disco”. Portanto, entre as alternativas disponíveis, a correta é a letra D.}
\else
\fi
\end{question}
}
%
{
\def\AMCbeginQuestion#1#2{\par\noindent{\bf (POSCOMP 2019 - Questão 29)}}
\begin{question}{} Em um sistema de arquivos, é necessário gerenciar blocos de discos que estão livres, para que possam ser alocados a arquivos e/ou diretórios de que necessitem. Analise as técnicas abaixo:
	
\begin{enumerate}[label=\Roman*.]
\item Vetor de bits, no qual cada bloco é representado por um bit.
\item Lista encadeada, na qual um ponteiro mantido em memória corresponde ao endereço do primeiro
bloco livre, e cada bloco contém um apontador para o próximo bloco livre.
\item Lista de contadores, em que cada elemento da lista contém um endereço de bloco e um contador
representando o número de blocos contíguos a este que estão livres.
\end{enumerate}

Quais implementam a gerência de espaço livre em disco?
\begin{multicols}{3}\columnseprule=0.1pt
\begin{choices}
\correctchoice{I, II e III.}
\wrongchoice{Apenas I.}
\wrongchoice{Apenas II.}
\wrongchoice{Apenas III.}
\wrongchoice{Apenas I e II.} 
\end{choices}
\end{multicols}
\end{question}
}
%
%
%
{
\def\AMCbeginQuestion#1#2{\par\noindent{\bf (POSCOMP 2019 – Questão 50)}}
\begin{question}{}  Uma partição de disco rígido é formatada com um sistema de arquivos com base em
alocação encadeada, usando tabela de alocação. A formatação da partição considerou um tamanho
de bloco (cluster) de 4096 bytes. Um programa executado nesse computador cria um arquivo nessa
partição e grava os dados do arquivo de acordo com o trecho de código abaixo. Como resultado, o
arquivo criado contém \underline{\hspace{1cm}} bytes e ocupa \underline{\hspace{1cm}} bytes na área de dados do sistema de arquivos.
\begin{lstlisting}[language=C]
...
05 char dado;
...
25 for(i=0;i<1024;i++)
27 write(fd, &dado, 1);
28 close(fd);
\end{lstlisting}

Assinale a alternativa que preenche, correta e respectivamente, as lacunas do trecho acima.
\begin{multicols}{3}\columnseprule=0.1pt
\begin{choices}
\correctchoice{1024 – 4096}
\wrongchoice{1024 – 1024}
\wrongchoice{1025 – 1025}
\wrongchoice{1025 – 4096}
\wrongchoice{4096 – 4096} 
\end{choices}
\end{multicols}
\end{question}
}
%
%
%
{
\def\AMCbeginQuestion#1#2{\par\noindent{\bf (POSCOMP 2018 - Questão 34)}}
\begin{question}{} Quanto aos métodos de alocação de espaço em disco para arquivos, analise as afirmações abaixo e assinale V, se verdadeiras, ou F, se falsas.
	
\begin{itemize}[label={( )}, leftmargin=1.5cm]
\item Na alocação contígua, a alocação de espaço em disco para novos arquivos pode ser dificultada pelo problema de fragmentação externa.
\item A alocação interligada provê acesso eficiente tanto a arquivos de acesso sequencial quanto de acesso direto.
\item Na alocação indexada, cada arquivo possui um bloco de índice. Para arquivos grandes, são necessários mecanismos que vinculem diferentes blocos de índices, como índices multinível, por exemplo.
\end{itemize}

A ordem correta de preenchimento dos parênteses, de cima para baixo, é:
\begin{multicols}{5}\columnseprule=0.1pt
\begin{choices}
\correctchoice{V – F – V.}
\wrongchoice{V – V – F.}
\wrongchoice{F – V – V.}
\wrongchoice{V – F – F.}
\wrongchoice{F – F – V.} 
\end{choices}
\end{multicols}
\end{question}
}
%
{
\def\AMCbeginQuestion#1#2{\par\noindent{\bf (POSCOMP 2017 - Questão 46)}}
\begin{question}{} Uma partição de disco rígido é formatada com um sistema de arquivos que utiliza
alocação encadeada baseada em tabela de alocação de arquivos (FAT). Após a formatação, a partição
possui setores de 512 bytes e tamanho de bloco (cluster) de 2048 bytes. Ao criar um arquivo nessa
partição, gravar 1 byte e fechá-lo, qual espaço esse arquivo ocupa na área de dados da partição?
	
\begin{multicols}{3}\columnseprule=0.1pt
\begin{choices}
\correctchoice{2048 bytes.}
\wrongchoice{1 byte.}
\wrongchoice{2 bytes.}
\wrongchoice{512 bytes.}
\wrongchoice{1 setor.}
\end{choices}
\end{multicols}
\end{question}
}
%
{
\def\AMCbeginQuestion#1#2{\par\noindent{\bf (POSCOMP 2016 - Questão 46})}
\begin{question}{} Um VSNT (Veículo Submarino Não Tripulado) é usado para monitoramento de
plataformas de petróleo marítimas. O VSNT tira uma foto a cada 1 minuto. O tamanho de cada arquivo
de foto é padronizado em 5 kB. As fotos são armazenadas em uma partição do disco rígido do VSNT,
a qual é formatada com sistema de arquivos FAT32 e tamanho de bloco (cluster) de 4 kB. O tempo
de missão do VSNT é de uma hora. Após o término de cada missão, as fotos são copiadas do VSNT
para um computador, que utiliza uma partição FAT32 formatada com clusters de 8 kB. Com base
nesse cenário, o espaço necessário no computador para armazenar todos os arquivos do VSNT em
uma missão é de:
\begin{tcolorbox}[
	colback=white,
	colframe=black,
	boxrule=0.8pt,
	arc=0pt,
	width=0.5\textwidth,
	left=8pt, right=8pt, top=6pt, bottom=6pt
	]
	\noindent kB: kilobyte \\
	1 kB = 1024 bytes
\end{tcolorbox}
\begin{multicols}{5}\columnseprule=0.1pt
\begin{choices}
\correctchoice{480 kB.}
\wrongchoice{600 kB.}
\wrongchoice{240 kB.}
\wrongchoice{300 kB.}
\wrongchoice{360 kB.} 
\end{choices}
\end{multicols}
\end{question}
}
{
\def\AMCbeginQuestion#1#2{\par\noindent{\bf (POSCOMP 2015 - Questão 46)}}
\begin{question}{} Considere um cenário de um sistema operacional que implementa um sistema de arquivos com método de alocação de espaço em disco baseado na alocação encadeada, a exemplo do popular sistema de arquivos FAT (file allocation table). Em um disco rígido com tamanho de setor igual a 512 bytes, criou-se uma partição e a formatou com esse sistema de arquivos usando 2048 bytes para o tamanho de blocos (clusters). Durante a escrita de dados em diferentes arquivos nessa partição, foi criado o arquivo ARQ.DAT que, após ter todos os seus dados armazenados, totalizou 1024 bytes de tamanho. Nesse cenário, o arquivo ARQ.DAT
	
\begin{choices}
\correctchoice{possui tamanho que não permite que seu conteúdo esteja fragmentado no disco.}
\wrongchoice{não possui seu conteúdo fragmentado no disco, pois o sistema de arquivos em uso não permite a fragmentação.}
\wrongchoice{pode ter seu conteúdo fragmentado no disco, pois já existiam outros arquivos no disco durante a sua criação e gravação, e o
sistema de arquivos em uso permite a fragmentação.}
\wrongchoice{pode ter seu conteúdo fragmentado no disco, pois seus dados foram armazenados concomitantemente com o armazenamento
de dados de outros arquivos, e o sistema de arquivos em uso permite a fragmentação.
}
\wrongchoice{pode ter seu conteúdo fragmentado no disco, pois seus dados ocupam, no mínimo, dois setores e o sistema de arquivos em uso
permite a fragmentação.
} 
\end{choices}
\end{question}
}
%
%
{
\def\AMCbeginQuestion#1#2{\par\noindent{\bf (POSCOMP 2013 - Questão 22)}}
\begin{question}{} Sobre arquivos e a alocação contígua em disco, considere as afirmativas a seguir.

\begin{enumerate}[label=\Roman*.]
\item Exige que se armazene o número do primeiro bloco do arquivo.
\item Fornece um desempenho excelente em operações de leitura.
\item Melhora o desempenho de acesso aleatório aos arquivos.
\item Minimiza a fragmentação de disco.
\end{enumerate}

Assinale a alternativa correta.
\begin{multicols}{2}\columnseprule=0.1pt
\begin{choices}
\correctchoice{Somente as afirmativas I, II e III são corretas.}
\wrongchoice{Somente as afirmativas I e II são corretas.}
\wrongchoice{Somente as afirmativas I e IV são corretas.}
\wrongchoice{Somente as afirmativas III e IV são corretas.}
\wrongchoice{Somente as afirmativas II, III e IV são corretas.} 
\end{choices}
\end{multicols}
\end{question}
}
%\input{EntradaSaida}
}


%##################################################################%
% Início da impressão de cada prova
\onecopy{1}{
\disciplina{PCS3746}{Sistemas Operacionais}
\avaliacao{Exercícios Gerenciamento de Processos}
%\cabecalho{}

\noindent
\AMCnumero{1} % para os testes começarem no 1
\vspace{0.5cm}
\shufflegroup{testes} % Para facilitar a revisão
\insertgroup{testes}
}
\end{document}